% 02_setup.tex -- Section 2, Setup (compressed; assets table + full defect text in Appendix A)
% Red lines unchanged. Notation kept in main text (Welch band + SNR definitions are load-bearing).

\section{Setup: Audited Claims, Assets, and Declared Defects}
\label{sec:setup}

We audit the three instruments as they are \emph{used} (at-a-glance map: Figure~\ref{fig:card}). Terminology (look-up table:
Table~\ref{tab:terms}, Appendix~\ref{app:setup}): a \emph{decision
instrument} is any cheap measurement whose reading stands in for an expensive one when
selecting data recipes or certifying ablations; an instrument's \emph{home setting} is the
configuration in which its advertised performance was established; and its \emph{usage
license} is the set of conditions under which its readings remain trustworthy (stated per
instrument in Sections~\ref{sec:t1}--\ref{sec:t3}). Further terms as they appear: an
\emph{arm} is one random-source condition (init-only, order-only, or bundled); a \emph{cell}
one task$\times$scale$\times$recipe point; a \emph{readout} the metric family (accuracy vs
bits-per-byte). \textbf{The checkpoint-noise proxy}: Signal \&
Noise \citep{signalandnoise2025} reports that the relative standard deviation of a run's final
checkpoints correlates with run-to-run seed noise at $R^2 = 0.82$--$0.95$ (i.e., $R \ge 0.9$),
and proposes the last-$n$-checkpoint spread as a cheap substitute for seed retraining.
\textbf{The init-only noise floor}: \citet{tencent2026} judge all 20 of their architecture claims
against a single 3-seed baseline noise floor, ``the seed-level distribution as the significance
threshold'' (their \S3.3, recommendation R5). The floor's source semantics are stated
inconsistently (\S3.3 varies ``both initialization and data-shuffle RNG'' per seed; their
appendix holds ``the shuffle seed fixed across all runs'' --- i.e.\ init-only as implemented), but either way the floor stands in
for the noise of \emph{all} random sources combined;
the auditable question is whether a single-source floor can play that role. Their per-seed
evaluations are not public as of August 2026; we audit the implemented init-only reading
(a bundled reading shifts Section~\ref{sec:t2}'s question, not its arithmetic).
\textbf{Small-scale
recipe ranking}: DataDecide \citep{datadecide2025} reports that the
ranking of 25 data recipes at a single 150M proxy scale predicts the better recipe at 1B with ${\sim}80\%$ pairwise accuracy (our reproduction: $80.2\%$ under their single-seed protocol, $82.7\%$ under 3-seed means; Section~\ref{sec:t3}); we audit the residual ${\sim}20\%$.

\paragraph{Assets and declared defects.}
All evidence comes from three public assets released with per-seed resolution:
\textbf{DataDecide} (14 sizes, 4M--1B, $\times$ 25 recipes $\times$ 3 seeds; OLMES suite
\citep{gu2025olmes} and
11-domain perplexities \citep{magnusson2024paloma}; $100$ tokens per parameter, ${\sim}5\times$Chinchilla
\citep{hoffmann2022chinchilla});
\textbf{Signal \& Noise \texttt{random\_seeds}} (20 1B runs of OLMo-2-style models
\citep{olmo2025}: 10 init-only, 9 data-order-only, 1 dense); and \textbf{PolyPythias} (50
Pythia \citep{biderman2023pythia} runs, 5 sizes $\times$ 10 seeds).
Table~\ref{tab:assets} (Appendix~\ref{app:setup}) inventories them with the defects that
constrain our methods, declared before any result.

\paragraph{Notation and estimands.}
\label{sec:setup-notation}
Throughout, $\sigma$ denotes a standard deviation of final scores across runs;
\emph{relative SD} divides by the arm mean. The proxy's \emph{calibration ratio} is
$y/x$ (true across-seed relative SD over the checkpoint-spread proxy). For
pairwise decisions we use pair-specific 95\% Welch
intervals for the difference of two 3-seed means: a pair $(i, j)$ sits inside its noise band
when $|\bar y_i - \bar y_j| < t_{0.975,\,\hat\nu_{ij}}\,\sqrt{\hat\sigma_i^2/3 +
\hat\sigma_j^2/3}$, with $\hat\nu_{ij}$ the Welch--Satterthwaite degrees of freedom. Bands within a cell are not
independent: the 300 bands of a 25-recipe cell share 25 recipe-level variance estimates, so we
read band-membership counts as descriptive tallies, not as independent tests.

\paragraph{The separability SNR.}
The
headline quantity of
Section~\ref{sec:t3} is the \emph{separability SNR}: signal $=$ the SD of the 25 recipe
means at a given proxy scale; noise $=$ the median within-recipe 3-seed SD. Because each
recipe mean itself carries seed variance $\hat\sigma_i^2/3$, the observed signal counts noise
as dispersion; we therefore also report the \emph{deconvolved} signal
$\sqrt{\max(S^2 - \overline{\hat\sigma^2}/3,\,0)}$, whose SNR uses the RMS noise
$\sqrt{\overline{\hat\sigma^2}}$ (diagnostics, robustness variants, and bootstrap intervals:
Appendix~\ref{app:t3}, Table~\ref{tab:snr-ci}). The separability SNR is target-free: it measures
separability of the proxy's readings, not correctness of its ranking (the transfer metrics carry
correctness).
